\documentclass[10pt,twocolumn,a4paper]{article}

\usepackage[margin=1.9cm,columnsep=0.75cm]{geometry}
\usepackage[utf8]{inputenc}
\usepackage[T1]{fontenc}
\usepackage{lmodern}
\usepackage{microtype}
\usepackage{amsmath,amssymb,amsthm,mathtools}
\usepackage{thmtools}
\usepackage{graphicx}
\usepackage{xcolor}
\usepackage{booktabs}
\usepackage{array}
\usepackage{enumitem}
\usepackage[numbers,sort&compress]{natbib}
\usepackage[colorlinks=true,linkcolor=blue!50!black,citecolor=blue!50!black,urlcolor=blue!50!black]{hyperref}
\usepackage[capitalise,noabbrev]{cleveref}
\usepackage{caption}
\captionsetup{font=small,labelfont=bf}

\declaretheorem[numberwithin=section]{theorem}
\declaretheorem[sibling=theorem]{proposition}
\declaretheorem[sibling=theorem]{lemma}
\declaretheorem[sibling=theorem]{corollary}
\declaretheorem[sibling=theorem,style=definition]{definition}
\declaretheorem[style=definition]{postulate}
\declaretheorem[sibling=theorem,style=definition]{example}
\declaretheorem[sibling=theorem,style=remark]{remark}

\newcommand{\N}{\mathbb{N}}
\newcommand{\Z}{\mathbb{Z}}
\newcommand{\R}{\mathbb{R}}
\newcommand{\kB}{k_{\mathrm{B}}}
\newcommand{\med}{m}
\newcommand{\sep}{\sigma}
\newcommand{\cut}{\partial}

\title{Positive Occupation Propagation:\\[4pt]
\large Direction, Time, Entropy and Individuation from Counted Commitment}
\author{%
Kundai Farai Sachikonye\\[0.3em]
\small AIMe Registry for Artificial Intelligence\\
\small Experimental Bioinformatics, Maximus-von-Imhof-Forum 3\\
\small 85354 Freising, Germany\\
\small \texttt{kundai.sachikonye@bitspark.com}}
\date{}

\begin{document}
\twocolumn[
\begin{@twocolumnfalse}
\maketitle
\begin{abstract}
\noindent
We examine what follows from one elementary arithmetic fact: a sum of
infinitely many positive integers has partial sums that increase by at least
one at every term and are therefore unbounded. We first state the fact
precisely and show that integrality, rather than positivity, carries its
content: a sequence of positive terms has a divergent sum along every infinite
selection of its terms if and only if the terms are bounded below by a
positive number, and positive real terms without such a bound can sum to a
finite total. We then model the physical act of recording---a
\emph{commitment}---as a term of such a sum, distinguish it from
\emph{occupation}, the cost-free presence of a system in a state, and define
\emph{positive occupation propagation} as the passage of occupation along a
contact that is accompanied by a commitment. From these definitions we derive,
under stated assumptions: a total order on recorded events (direction); a
time defined as a ratio of commitment counts; a register entropy that cannot
decrease; and an individuation measure, the minimum cut separating an item
from a medium in a contact graph, which is bounded below by one when contact
weights are counted and which cannot decrease as the graph grows. We show that
the corresponding statement for real-valued contact weights is false, and give
the counterexample. Two classical problems are re-examined in this setting.
In the dichotomy of Zeno the distances form a convergent measured series while
the commitments needed to distinguish the midpoints form a divergent counted
series; the paradox's premises refer to the second. In the reversibility
objection of Loschmidt, configurations of a finite reversible system recur and
can be reversed, whereas the count of commitments made while doing so
increases throughout. Finally we write a sequence of successive refinements,
each removing a fraction of what remains, as a \emph{ladder}, and show that its
logarithmic gap is a sum of positive terms that diverges exactly when the
fractions are bounded below. Each statement is accompanied by a computational
check: partial sums over $10^4$ terms and $10^3$ trials, separation costs in
$6{,}400$ growth steps of random contact graphs, recurrence and reversal in a
Kac ring and in the cat map on lattices up to $400\times400$, and exact
rational ladders. The scope is deliberately narrow: the results concern
records, orderings and cuts, and do not determine any dynamical law.
\end{abstract}
\vspace{1.2em}
\end{@twocolumnfalse}
]


%==============================================================================
\section{Introduction}
\label{sec:intro}
%==============================================================================

The equations of classical and quantum mechanics are invariant under reversal
of the time parameter, yet records of what has happened only accumulate. A
laboratory notebook gains pages; a detector's count register is incremented;
a photographic plate darkens. The tension between the two observations is old
and has been discussed under several names---the reversibility objection of
Loschmidt to Boltzmann's $H$-theorem~\citep{loschmidt1876,boltzmann1877},
the problem of the arrow of time~\citep{lebowitz1993,price1996,albert2000},
and, in a different register, Zeno's arguments against the composition of
motion from points~\citep{grunbaum1967,salmon2001}. These discussions share a
feature: each turns on what it means to \emph{distinguish} one state, one
instant or one point from another, and on what such distinctions cost.

This article isolates one elementary fact and traces, as carefully as we can,
what it does and does not imply for those questions. The fact is arithmetic:
\begin{quote}
\emph{a sum of infinitely many positive integers has partial sums that
increase by at least one at every term, and therefore exceed every bound.}
\end{quote}
Stated this way the fact is trivial. Its interest lies in which of its words
carry the weight. ``Positive'' alone does not: $\tfrac12+\tfrac14+\tfrac18+\cdots$
has positive terms and a finite sum. ``Integer'' does, because it supplies a
uniform lower bound of one. ``Infinitely many'' supplies unboundedness, and
``each term an integer'' supplies finiteness of each term. We make these
roles precise in \cref{sec:sums}, where the central statement
(\cref{thm:selection}) is that a sequence of positive terms has a divergent
sum along every infinite selection of its terms if and only if its terms are
bounded below by a positive number.

The physical content enters through one modelling decision, set out in
\cref{sec:model}: the act of making a record---which we call a
\emph{commitment}---is counted, in the sense that each commitment costs at
least a fixed positive amount $\theta$. We keep this separate from
\emph{occupation}: a system is in some state at every moment, and being there
is not a record and has no cost in this accounting. The term \emph{positive
occupation propagation} denotes the passage of occupation from one item to
another along a contact, accompanied by a commitment; it is the operation
whose consequences the rest of the article follows.

\paragraph{What is derived.}
Under the definitions of \cref{sec:model} we obtain:
\begin{enumerate}[label=(\roman*),nosep]
\item a strict total order on commitments, which we call direction
      (\cref{sec:direction});
\item a precise account of which series the dichotomy and the arrow refer to
      (\cref{sec:zeno});
\item a time defined as a ratio of commitment counts, of the kind used in
      the definition of the SI second (\cref{sec:time});
\item a register entropy that increases by at least $\kB\ln b$ per commitment
      (\cref{sec:entropy});
\item an individuation measure on contact graphs that is at least one under
      counted weights and does not decrease when the graph grows, together
      with a counterexample showing that the second property fails to
      protect a lower bound under real weights (\cref{sec:individuation});
\item the joint monotonicity of these quantities under positive occupation
      propagation (\cref{sec:up});
\item the recurrence of configurations and non-recurrence of counts in
      finite reversible systems, and the reading of Loschmidt's objection
      this suggests (\cref{sec:bounded});
\item a ladder calculus for successive refinement and the condition under
      which a ladder's remaining gap tends to zero (\cref{sec:ladders}).
\end{enumerate}
\Cref{sec:devices} relates each derived quantity to the instrument that
reads it, \cref{sec:validation} reports the computational checks, and
\cref{sec:limits} states the limitations.

\paragraph{What is not claimed.}
Nothing here derives a dynamical law, a physical constant, or the numerical
value of the floor $\theta$; where a physical value is used (the Landauer
bound in \cref{sec:entropy}) it is imported from established physics and
marked as such. The treatment of Zeno's paradoxes does not replace the
analytic resolution through convergent series and limits; it identifies, for
each premise, which kind of series the premise concerns. The treatment of
Loschmidt's objection does not dispute the statistical explanation of the
second law; it separates two quantities that the objection is sometimes taken
to treat as one.

\paragraph{Notation.}
$\N=\{1,2,3,\dots\}$; $\Z_{\ge 1}$ is the same set when emphasising
integrality. For a sequence $(c_i)_{i\ge1}$ we write $S_n=\sum_{i=1}^n c_i$.
Graphs are finite, undirected and simple unless stated; $\cut S$ denotes the
set of edges with exactly one endpoint in $S$, and $w(\cut S)$ its total
weight.

%==============================================================================
\section{Counted and measured sums}
\label{sec:sums}
%==============================================================================

\begin{definition}[Cost sequences]
\label{def:costs}
A \emph{cost sequence} is a sequence $(c_i)_{i\ge1}$ of positive real numbers.
It is \emph{counted} if every $c_i\in\Z_{\ge1}$; \emph{floored} if there is
$\theta>0$ with $c_i\ge\theta$ for all $i$; and \emph{measured} if only
$c_i>0$ is assumed.
\end{definition}

\begin{lemma}[Partial sums of a counted sequence]
\label{lem:counted}
If $(c_i)$ is counted then $S_{n+1}-S_n\ge1$ and $S_n\ge n$ for all $n$. In
particular $(S_n)$ is strictly increasing and unbounded.
\end{lemma}

\begin{proof}
$S_{n+1}-S_n=c_{n+1}\ge1$. Summing the first $n$ increments from $S_0=0$
gives $S_n\ge n$, which exceeds any bound for large $n$.
\end{proof}

\begin{proposition}[Floored reduces to counted]
\label{prop:floored}
If $(c_i)$ is floored with floor $\theta$, then $k_i:=\lfloor c_i/\theta\rfloor$
is counted and $S_n\ge\theta\sum_{i\le n}k_i\ge\theta n$.
\end{proposition}

\begin{proof}
$c_i/\theta\ge1$ gives $k_i\ge1$, and $c_i\ge\theta k_i$.
\end{proof}

The choice of integers is therefore a choice of unit: any uniform positive
lower bound can be rescaled to one. What cannot be rescaled away is the
\emph{absence} of such a bound.

\begin{proposition}[Positivity is not a floor]
\label{prop:positivity}
There are measured sequences with
\begin{enumerate}[label=(\alph*),nosep]
\item a finite sum, for example $c_i=2^{-i}$ with $\sum c_i=1$;
\item a divergent sum but $\inf c_i=0$, for example $c_i=1/i$.
\end{enumerate}
\end{proposition}

\begin{proof}
Both are standard: the geometric series converges and the harmonic series
diverges~\citep{knopp1990}.
\end{proof}

The harmonic example shows that divergence of the whole sum does not imply a
floor. The following statement identifies exactly what a floor buys.

\begin{theorem}[Divergence along every selection]
\label{thm:selection}
For a measured sequence $(c_i)$ the following are equivalent:
\begin{enumerate}[label=(\roman*),nosep]
\item $\inf_i c_i>0$;
\item for every infinite set $I\subseteq\N$, $\sum_{i\in I}c_i=\infty$;
\item for every infinite set $I\subseteq\N$, the partial sums of
      $(c_i)_{i\in I}$ in increasing order of $i$ grow at least linearly in
      the number of terms.
\end{enumerate}
\end{theorem}

\begin{proof}
(i)$\Rightarrow$(iii): with $\theta=\inf c_i>0$, any $n$ terms sum to at least
$n\theta$. (iii)$\Rightarrow$(ii) is immediate. (ii)$\Rightarrow$(i): suppose
$\inf c_i=0$. Then for every $k$ there are infinitely many $i$ with
$c_i<2^{-k}$, so we may choose $i_1<i_2<\cdots$ with $c_{i_k}<2^{-k}$. The
selection $I=\{i_k\}$ has $\sum_{i\in I}c_i<1$, contradicting (ii).
\end{proof}

\Cref{thm:selection} is the precise form of the statement in the
introduction. A process whose individual steps carry a uniform positive cost
accumulates cost without bound along \emph{any} infinite subsequence of its
steps; a process whose steps are merely positive does not, because one can
always select infinitely many of its steps whose total cost is finite. The
second possibility is exactly what Zeno's dichotomy exploits
(\cref{sec:zeno}).

\begin{remark}[Finiteness of each term]
\label{rem:finite-terms}
Each term of a counted sequence is a finite integer. In what follows this
will correspond to the requirement that every single commitment, and every
contact weight, is finite. Infinitely many terms supply an unbounded total;
each term separately is bounded.
\end{remark}

\Cref{fig:sums} shows the three regimes numerically: counted partial sums
lying on or above the line $S_n=n$ in every one of $10^3$ trials of $10^4$
terms; the geometric series stalling at one; the harmonic series diverging
slowly; and the running minimum of positive real terms $c_i=U_i/i$ falling
below $10^{-7}$, while the running minimum of the counted terms stays at one.

\begin{figure*}[t]
\centering
\includegraphics[width=\textwidth]{figures/f1_partial_sums.pdf}
\caption{\textbf{Counted and measured sums.}
(\textbf{A})~Partial sums of counted costs $c_i\in\{1,\dots,5\}$ (median and
95\% band over $10^3$ sequences of $10^4$ terms) lie above $S_n=n$; no
violation occurred. The harmonic and geometric real sequences are shown for
comparison.
(\textbf{B})~Running minimum of the terms. For positive reals $c_i=U_i/i$,
$U_i$ uniform on $(0,1)$, the median running minimum after $10^4$ terms is
$1.4\times10^{-8}$; for counted terms it is one.
(\textbf{C})~The dichotomy (\cref{sec:zeno}): the number of committed
midpoints needed to come within $\varepsilon$ of the goal,
$k(\varepsilon)=\lceil\log_2(D/\varepsilon)\rceil$, is finite for every
$\varepsilon>0$ and unbounded as $\varepsilon\to0$, while the distance covered
(inset) converges to $D$.}
\label{fig:sums}
\end{figure*}

%==============================================================================
\section{Occupation, commitment and propagation}
\label{sec:model}
%==============================================================================

We now state the model. Its ingredients are a set of states a system can
occupy, a register in which records are made, and a contact relation along
which occupation can pass.

\begin{definition}[Occupation]
\label{def:occupation}
Let $X$ be a set of states. An \emph{occupation history} is a map
$x:\mathcal{T}\to X$ from an ordered index set $\mathcal{T}$; $x(\tau)$ is the
state occupied at $\tau$. Occupation is not itself a record: no cost is
attached to a system's being in a state.
\end{definition}

\begin{definition}[Register and commitment]
\label{def:register}
A \emph{register} over a finite alphabet $A$, $|A|=b\ge2$, is a finite word
$R\in A^{*}$. A \emph{commitment} is an event that appends one letter to a
register. The \emph{commitment count} $M$ of a register is the number of
commitments made to it.
\end{definition}

The count $M$ counts events, not the current length of the word: if a
register can be overwritten, each overwrite is itself an event and is
counted. With this reading $M$ is non-decreasing by construction, and the
content of the model lies entirely in the following assumption about cost.

\begin{postulate}[Counted commitment]
\label{post:counted}
There is a constant $\theta>0$ such that every commitment costs at least
$\theta$, and no operation that is not a commitment changes any register.
\end{postulate}

Two remarks fix the status of \cref{post:counted}. First, it is an
assumption, and it is the only physical assumption of the article; everything
else is definition and proof. Second, it has a well-established physical
counterpart: resetting one bit of a physical memory in contact with a heat
bath at temperature $T$ dissipates at least $\kB T\ln2$ of
heat~\citep{landauer1961,bennett1982}, a bound that has been approached in
experiment~\citep{berut2012}. A register of bounded physical size that keeps
recording must reuse its cells, and each reuse is subject to that bound. We
use $\theta$ abstractly and do not need its value.

\begin{definition}[Contact graph with medium]
\label{def:contact}
A \emph{contact graph} is a finite connected weighted graph $G=(V,E,w)$ with
$w:E\to(0,\infty)$ and a distinguished vertex $\med\in V$, the \emph{medium}.
Its weights are \emph{counted} if $w(E)\subseteq\Z_{\ge1}$.
\end{definition}

The medium stands for whatever an item is distinguished against: the solvent
around a molecule, the background of an image, the rest of a sample. We do
not require it to be adjacent to every vertex, only that the graph be
connected.

\begin{definition}[Propagation and positive occupation propagation]
\label{def:pop}
A \emph{propagation step} moves occupation from a vertex $u$ to a neighbour
$v$ along an edge $\{u,v\}\in E$. It is a \emph{positive occupation
propagation} (POP) step if it is accompanied by a commitment that records it.
A POP step that creates a new contact---adds an edge, or adds a vertex and
its edges---is \emph{constructive}.
\end{definition}

\begin{remark}[Probe and commit]
\label{rem:probe}
A propagation step without a commitment is a \emph{probe}: occupation moves
but nothing is recorded. Probes are unrestricted in this accounting. All the
monotone quantities derived below are functions of commitments only;
configurations reached by probes can repeat (\cref{sec:bounded}).
\end{remark}

%==============================================================================
\section{Direction}
\label{sec:direction}
%==============================================================================

\begin{theorem}[Commitments are totally ordered]
\label{thm:order}
Let $e_1,e_2,\dots$ be the commitments made to a register, indexed in the
order in which they are made, and let $M(e_k)=k$ be the count after $e_k$.
Then $M$ is strictly increasing along the sequence, and the cumulative cost
$C_k$ satisfies $C_k\ge k\theta$ and $C_{k+1}-C_k\ge\theta$.
\end{theorem}

\begin{proof}
Strict increase of $M$ is the definition of the index. The cost statements
follow from \cref{post:counted} and \cref{lem:counted} applied to the
floored sequence of individual costs, via \cref{prop:floored}.
\end{proof}

We call the order given by $M$ the \emph{direction} of the register. Three
points deserve emphasis. (a) Direction is a property of commitments, not of
occupation: \cref{thm:order} says nothing about the order of states
occupied. (b) No commitment lowers $M$ or $C$, because no term of a counted
sum is zero or negative. An event that ``undoes'' a record---erasing a cell,
say---is itself a commitment and raises both. (c) By \cref{thm:selection},
every infinite set of commitments, however it is selected, carries unbounded
cost; there is no infinite collection of records that costs a finite amount.

\begin{corollary}[No finite infinity of records]
\label{cor:no-zeno-records}
A system with finite total budget $B$ makes at most $\lfloor B/\theta\rfloor$
commitments.
\end{corollary}

%==============================================================================
\section{The dichotomy and the arrow}
\label{sec:zeno}
%==============================================================================

Zeno's dichotomy, in one common form, argues that a runner cannot reach the
end of a course of length $D$: before reaching the end the runner must reach
the midpoint, before that the quarter point, and so on, so that infinitely
many tasks must be completed~\citep{grunbaum1967,salmon2001}. The standard
reply observes that the distances $D/2,D/4,\dots$ and the corresponding times
form convergent series, so that infinitely many sub-intervals occupy a finite
distance and a finite time. We do not dispute the reply. We ask instead which
kind of sum each premise of the argument refers to.

\begin{definition}[The dichotomy, as occupation and as commitment]
\label{def:dichotomy}
Let the course be $[0,D]$ and the dichotomy points $p_k=D(1-2^{-k})$,
$k\ge1$. The runner \emph{occupies} $p_k$ if its occupation history passes
through $p_k$; the runner (or any observer) \emph{commits} $p_k$ if a
commitment records that $p_k$ was reached.
\end{definition}

\begin{proposition}[Two series in one argument]
\label{prop:two-series}
\begin{enumerate}[label=(\alph*),nosep]
\item The distances between successive dichotomy points form a measured
      sequence with sum $D$.
\item Committing the first $k$ dichotomy points costs at least $k\theta$;
      committing all of them is impossible with any finite budget.
\item To commit a point within $\varepsilon>0$ of the goal along the
      dichotomy requires exactly
      $k(\varepsilon)=\lceil\log_2(D/\varepsilon)\rceil$ commitments, which is
      finite for every $\varepsilon>0$.
\end{enumerate}
\end{proposition}

\begin{proof}
(a) is the geometric series. (b) is \cref{thm:order} and
\cref{cor:no-zeno-records}. (c): $D-p_k=D2^{-k}\le\varepsilon$ iff
$k\ge\log_2(D/\varepsilon)$.
\end{proof}

The argument's conclusion---that the runner must \emph{complete} infinitely
many \emph{tasks}---treats reaching each point as an act that must be done
and finished before the next, that is, as a commitment. The analytic reply
treats reaching each point as occupation, measured by a convergent series of
lengths. \Cref{prop:two-series} shows that both readings are internally
consistent and that they disagree only because they concern different sums:
occupation of all dichotomy points is completed in finite distance and time,
and commitment of all of them is never completed. A runner that only
occupies arrives; an observer that commits each point never finishes
committing, but for any resolution $\varepsilon$ finishes the commitments
needed to resolve the arrival to that resolution, after finitely many
(\cref{fig:sums}C; $k=50$ commitments for $\varepsilon/D=10^{-15}$).

The arrow paradox---that a flying arrow occupies at each instant a space equal
to itself and is therefore at rest at each instant---can be placed in the
same scheme. A committed instant records a position. A record of position is
not a record of motion, and no finite set of position records contains a
velocity; a velocity is a relation between records, or a property of the
occupation history between them. The paradox's inference ``at rest at every
instant, hence at rest'' composes properties of records into a property of
the occupation history, and the composition is not valid in either
direction. We do not regard this as more than a restatement of the
well-known point that motion is a relation between positions at different
times~\citep{salmon2001}; it is included because it uses only the
occupation--commitment distinction already introduced.

%==============================================================================
\section{Time as a count}
\label{sec:time}
%==============================================================================

\begin{definition}[Clock and recorded time]
\label{def:clock}
A \emph{clock} is a reference process that makes commitments to its own
register at a fixed number $f$ per unit cycle. The \emph{recorded time}
between two commitments $e,e'$ of the clock is $t(e,e')=(M(e')-M(e))/f$.
\end{definition}

This is not an exotic definition. The SI second is defined by fixing the
numerical value of the caesium-133 ground-state hyperfine transition
frequency at $9\,192\,631\,770$~Hz~\citep{bipm2019}; an interval is measured
as one second when that many cycles have been counted. Recorded time in the
sense of \cref{def:clock} is a count divided by a fixed count.

\begin{proposition}[Properties of recorded time]
\label{prop:time}
Recorded time is additive, $t(e,e'')=t(e,e')+t(e',e'')$ for $e\preceq
e'\preceq e''$; non-negative; zero only between a commitment and itself; and
unbounded along any infinite sequence of clock commitments.
\end{proposition}

\begin{proof}
Additivity is additivity of differences. Non-negativity and the zero case
follow from \cref{thm:order}. Unboundedness is \cref{lem:counted}.
\end{proof}

\begin{remark}
\label{rem:parameter}
Recorded time is not assumed to be the parameter appearing in a dynamical
law. The two are related by calibration: a dynamical model is accepted if its
parameter, mapped through the dynamics of a reference process, reproduces
the counts. The asymmetry of recorded time is a property of counts and does
not require any asymmetry in the parameter.
\end{remark}

%==============================================================================
\section{Register entropy}
\label{sec:entropy}
%==============================================================================

\begin{definition}[Register entropy]
\label{def:regent}
A register that has received $M$ commitments over an alphabet of size $b$
distinguishes $b^M$ possible histories of commitments. Its \emph{register
entropy} is $S_{\mathrm{reg}}=\kB\ln b^M=\kB M\ln b$.
\end{definition}

The definition is Boltzmann's $S=\kB\ln W$ applied to the set of
distinguishable commitment histories, and the quantity $M\log_2 b$ is the
Shannon information capacity of the record in bits~\citep{shannon1948}.

\begin{proposition}[Register entropy does not decrease]
\label{prop:regent}
$S_{\mathrm{reg}}$ is non-decreasing in the commitment order and increases by
exactly $\kB\ln b$ at each commitment.
\end{proposition}

\begin{proof}
Immediate from \cref{def:regent} and \cref{thm:order}.
\end{proof}

\begin{corollary}[Three readings of one count]
\label{cor:three}
For a register written at the rate of a clock with frequency $f$,
\begin{equation}
M \;=\; f\,t \;=\; \frac{S_{\mathrm{reg}}}{\kB\ln b}.
\label{eq:three}
\end{equation}
Count, recorded time and register entropy are proportional; none can
increase while another decreases.
\end{corollary}

\begin{remark}[Register entropy and configurational entropy]
\label{rem:two-entropies}
Register entropy must not be confused with the Boltzmann entropy of a
system's macrostate, $S_{\mathrm{conf}}=\kB\ln|\Gamma_{\mathrm{macro}}|$. The
latter is a function of the configuration occupied and can decrease, rarely
when unprepared and at a cost when prepared. The former is a function of the
record and cannot decrease. \Cref{sec:bounded} exhibits the two side by
side in a reversible system.
\end{remark}

%==============================================================================
\section{Individuation in contact graphs}
\label{sec:individuation}
%==============================================================================

We now ask how an item is distinguished from its medium. The measure we use
is the cheapest way of cutting the item off.

\begin{definition}[Separation cost]
\label{def:sep}
For a contact graph $(G,\med)$ and $v\in V\setminus\{\med\}$, the
\emph{separation cost} of $v$ is
\begin{equation}
\sep_G(v)\;=\;\min\bigl\{\,w(\cut S)\;:\;v\in S\subseteq V\setminus\{\med\}\,\bigr\}.
\label{eq:sep}
\end{equation}
A minimiser $S^{*}(v)$ is an \emph{identity region} of $v$.
\end{definition}

By the max-flow min-cut theorem~\citep{ford1956}, $\sep_G(v)$ equals the
value of a maximum flow from $v$ to $\med$ with edge capacities $w$, and is
computable in polynomial time.

\begin{theorem}[Floor of individuation]
\label{thm:floor}
If $G$ is connected and its weights are counted, then for every
$v\ne\med$,
\begin{equation}
\sep_G(v)\;\ge\;\lambda_G(v,\med)\;\ge\;1,
\end{equation}
where $\lambda_G(v,\med)$ is the maximum number of edge-disjoint paths from
$v$ to $\med$.
\end{theorem}

\begin{proof}
Every set $S$ with $v\in S$, $\med\notin S$ is crossed by each of
$\lambda_G(v,\med)$ edge-disjoint $v$--$\med$ paths at least once, so
$|\cut S|\ge\lambda_G(v,\med)$ (Menger~\citep{menger1927}). Counted weights
give $w(\cut S)\ge|\cut S|$. Connectivity gives $\lambda_G(v,\med)\ge1$.
\end{proof}

\begin{theorem}[Monotonicity under expansion]
\label{thm:monotone}
Let $G\subseteq G'$ be contact graphs with the same medium, $G'$ obtained
from $G$ by adding vertices and edges and leaving existing weights
unchanged. Then $\sep_{G'}(v)\ge\sep_G(v)$ for every $v\in V\setminus\{\med\}$.
\end{theorem}

\begin{proof}
Let $S'$ attain $\sep_{G'}(v)$ and put $S=S'\cap V$. Then $v\in S$ and
$\med\notin S$, and every edge of $G$ crossing $S$ is an edge of $G'$
crossing $S'$, since both endpoints lie in $V$ and their sides are
unchanged. Hence $\sep_G(v)\le w(\cut_G S)\le w(\cut_{G'}S')=\sep_{G'}(v)$.
\end{proof}

\Cref{thm:monotone} holds for any positive weights. What does \emph{not}
survive under real weights is any lower bound that is independent of the
graph:

\begin{proposition}[No invariant floor under measured weights]
\label{prop:no-floor}
There is an expanding sequence $G_0\subseteq G_1\subseteq\cdots$ of contact
graphs with positive real weights in which (a) the smallest weight incident to
a fixed vertex $v$ tends to zero, and (b) the separation costs of the added
vertices tend to zero.
\end{proposition}

\begin{proof}
Let $G_0$ have vertices $\{v,\med\}$ and one edge of weight one. Form
$G_k$ by adding a vertex $u_k$ with edges $\{u_k,v\}$ and $\{u_k,\med\}$ of
weight $1/k$ each. The smallest weight incident to $v$ in $G_k$ is $1/k$, and
$\sep_{G_k}(u_k)\le w(\cut\{u_k\})=2/k$.
\end{proof}

\begin{remark}[Why counted weights]
\label{rem:why-counted}
\Cref{prop:no-floor} is the graph form of \cref{prop:positivity}: positive
weights without a floor admit a sequence of contacts whose costs sum to a
finite amount. With counted weights the global minimum weight is one in every
graph of every size, so the floor of \cref{thm:floor} is the same however
large the graph becomes. A claim that a positive lower bound on individuation
persists as a graph grows therefore needs counted (or uniformly floored)
weights; positivity of each weight is not enough.
\end{remark}

\begin{proposition}[Identity is a region]
\label{prop:region}
There are contact graphs in which every identity region of some $v$ contains
a vertex other than $v$.
\end{proposition}

\begin{proof}
Take vertices $v,u,\med$ with $w(v,u)=10$, $w(v,\med)=3$, $w(u,\med)=1$. Then
$w(\cut\{v\})=13$ and $w(\cut\{v,u\})=4$. These are the only sets containing
$v$ and not $\med$, so $S^{*}(v)=\{v,u\}$ uniquely.
\end{proof}

\Cref{prop:region} says that what an item is, as measured by how it is cut
from its surroundings, is in general settled together with some of its
neighbours. Measurement can be modelled in the same terms: attaching a probe
to an item adds a vertex and edges, and by \cref{thm:monotone} a measurement
of this kind can never lower the separation cost of anything already in the
graph.

\Cref{fig:individuation} reports $40$ random growth sequences of $160$ steps
each, under counted weights (uniform on $\{1,\dots,5\}$) and real weights
(exponential, mean $1.5$). The separation cost of a fixed vertex never
decreased in any of the $12{,}800$ steps, under either weighting, as
\cref{thm:monotone} requires. The smallest weight in the graph stayed at one
under counted weights and fell to a median of $1.4\times10^{-3}$ (minimum
$1.1\times10^{-5}$) under real weights, as \cref{rem:why-counted} anticipates.
The smallest separation cost over all vertices of the final graphs was $5$
under counted weights and $0.86$ under real weights; identity regions larger
than a single vertex occurred for $0.7\%$ of vertices under counted weights
and $17\%$ under real weights.

\begin{figure*}[t]
\centering
\includegraphics[width=\textwidth]{figures/f2_individuation.pdf}
\caption{\textbf{Separation cost under expansion.}
Contact graphs grown by adding one vertex per step, joined to the medium and
to three existing vertices; with probability $0.15$ the fixed vertex $v$ also
gains the new vertex as neighbour.
(\textbf{A})~$\sep(v)$ for eight counted-weight sequences; no decrease
occurred in $40\times160$ steps under either weighting.
(\textbf{B})~Smallest edge weight in the graph (median, 10--90\% band).
(\textbf{C})~Separation costs of all vertices of the final graphs of ten
sequences per weighting; the dashed line marks one.}
\label{fig:individuation}
\end{figure*}

%==============================================================================
\section{Only up: direction and individuation together}
\label{sec:up}
%==============================================================================

Positive occupation propagation joins the two strands. Each POP step is a
commitment, so it advances the count; each constructive POP step adds
contact, so it cannot lower any separation cost.

\begin{theorem}[Joint monotonicity under positive occupation propagation]
\label{thm:up}
Let $(G_k,R_k)_{k\ge0}$ be the contact graph and register after the $k$-th
step of a sequence of POP steps, each either constructive or not, with
counted weights. Then for every vertex $v$ present from step $k_0$ on and
every $k\ge k_0$,
\begin{gather}
M_{k+1}=M_k+1,\qquad S_{\mathrm{reg},k+1}=S_{\mathrm{reg},k}+\kB\ln b,\\
\sep_{G_{k+1}}(v)\;\ge\;\sep_{G_k}(v)\;\ge\;1.
\end{gather}
\end{theorem}

\begin{proof}
The first two equalities are \cref{thm:order} and \cref{prop:regent}. A
non-constructive step leaves the graph unchanged; a constructive step
produces $G_{k+1}\supseteq G_k$, and \cref{thm:monotone} applies. The bound
by one is \cref{thm:floor}.
\end{proof}

The informal summary ``the only direction is up'' is, in this precise sense,
the statement that three quantities---the count, the register entropy, and
the separation cost of every item already present---are non-decreasing under
positive occupation propagation, the first two strictly. It is not a
statement about configurations, which are free to return (\cref{sec:bounded}),
nor about the separation cost of items added later, which start wherever
their contacts place them (but, under counted weights, not below one).

%==============================================================================
\section{Bounded parts, recurrence, and the reversibility objection}
\label{sec:bounded}
%==============================================================================

\Cref{rem:finite-terms} noted that each term of a counted sum is finite.
Here we use this to bound the number of distinguishable configurations of a
part, and then look at what reversible dynamics does on finitely many
configurations.

\begin{proposition}[Finite parts have finitely many distinguishable configurations]
\label{prop:finite}
If distinguishing between two alternatives costs at least $\theta$ and a part
has a finite budget $B$ for distinctions, then at most
$2^{\lfloor B/\theta\rfloor}$ of its configurations are mutually
distinguishable.
\end{proposition}

\begin{proof}
Distinguishing among $N$ alternatives by binary distinctions requires at least
$\lceil\log_2N\rceil$ of them, at cost at least $\theta$ each
(\cref{post:counted}); so $\theta\log_2N\le B$.
\end{proof}

We treat the resulting finite set of distinguishable configurations as the
state space of the part. On a finite set, an invertible deterministic
dynamics is a permutation, and the following is elementary.

\begin{proposition}[Recurrence on a finite set]
\label{prop:recurrence}
Let $F$ be a bijection of a finite set $X$. Every orbit of $F$ is periodic:
for each $x\in X$ there is $P\ge1$, $P\le|X|$, with $F^P(x)=x$.
\end{proposition}

For measure-preserving dynamics on a space of finite measure, the
corresponding statement is Poincaré's recurrence theorem~\citep{poincare1890},
with Kac's lemma giving the mean return time~\citep{kac1947}.

\begin{theorem}[Configurations recur; counts do not]
\label{thm:loschmidt}
Let $F$ be a bijection of a finite set $X$ and let a register receive one
commitment per application of $F$ or $F^{-1}$. Starting from $x_0$ with
$M=0$:
\begin{enumerate}[label=(\alph*),nosep]
\item after $P$ forward steps, where $P$ is the period of $x_0$, the
      configuration is $x_0$ and $M=P$;
\item after $T$ forward steps followed by $T$ steps of $F^{-1}$, the
      configuration is $x_0$ and $M=2T$.
\end{enumerate}
In both cases the configuration returns and the count does not.
\end{theorem}

\begin{proof}
(a) is \cref{prop:recurrence} with one commitment per step. (b):
$F^{-T}F^{T}=\mathrm{id}$, and $2T$ commitments have been made.
\end{proof}

\paragraph{Reading of the reversibility objection.}
Loschmidt objected that for every trajectory along which a macroscopic
entropy increases, the time-reversed trajectory, obtained by reversing all
velocities, is also a solution, and along it the entropy
decreases~\citep{loschmidt1876}. \Cref{thm:loschmidt} separates two things
the word ``reversal'' can denote.
\begin{itemize}[nosep]
\item \emph{Reversal of configuration} is permitted by reversible dynamics:
      the configuration returns to $x_0$ and the configurational entropy
      returns with it. Spontaneous return in a large system is atypical
      (its measure is small and Kac's mean return time is long); prepared
      return is possible but requires operations on the system, of which the
      spin echo is a laboratory example~\citep{hahn1950}.
\item \emph{Reversal of the count} is not an operation of the model. Every
      step, forward or backward, is a commitment, and $M$ increases by one
      at each.
\end{itemize}
On this reading, the statement ``reversal is possible but improbable'' and
the statement ``reversal of the record is not a well-defined operation'' are
both true, because they concern different quantities. We do not claim that
this resolves every question raised under the heading of the arrow of
time~\citep{lebowitz1993,albert2000}; only that the objection, as applied to
records, does not apply.

\paragraph{Numerical illustration.}
The Kac ring~\citep{kac1959,gottwald2009} is a minimal reversible system:
$N$ sites on a circle each carry a ball coloured black or white, a fixed set
of edges is marked, and at each step every ball moves one site clockwise,
changing colour when it crosses a marked edge. The dynamics is invertible
and every configuration has period dividing $2N$. With $N=400$ and $41$
marked edges, starting from all white, the coarse-grained entropy of the
colour fraction rises to its maximum within a few dozen steps, the
configuration returns exactly at step $800$, and at step $400$ every ball has
the opposite colour (\cref{fig:recurrence}A; an odd number of marks flips
each ball an odd number of times per circuit). Reversing the dynamics after
$T=150$ steps returns the initial configuration exactly at step $300$, with
the coarse-grained entropy falling back to zero, while the register received
$300$ commitments, one per step (\cref{fig:recurrence}B). As a second check of
\cref{prop:recurrence}, we computed the period of the cat map
$(x,y)\mapsto(x+y,\,x+2y)\bmod N$~\citep{arnold1968} on $N\times N$ lattices
for $N=2,\dots,400$: every period is finite and none exceeds $3N$, the bound
of Dyson and Falk~\citep{dyson1992} (\cref{fig:recurrence}C).

\begin{figure*}[t]
\centering
\includegraphics[width=\textwidth]{figures/f3_recurrence.pdf}
\caption{\textbf{Configurations recur; the register does not.}
(\textbf{A})~Kac ring, $N=400$, $41$ marked edges, initially all white:
coarse-grained entropy of the colour fraction (in units of $\ln2$) and
Hamming distance to the initial configuration. The configuration recurs at
step $800$ (dashed).
(\textbf{B})~Forward dynamics for $150$ steps, then the inverse dynamics
(dotted line): the configuration returns exactly at step $300$ while the
register length, one commitment per step, rises throughout.
(\textbf{C})~Period of every orbit of the cat map on an $N\times N$ lattice,
$N=2,\dots,400$, with the bound $3N$.}
\label{fig:recurrence}
\end{figure*}

%==============================================================================
\section{Ladders}
\label{sec:ladders}
%==============================================================================

Many processes refine in steps, each step removing a fraction of what
remains unresolved: successive purification, successive filtering, successive
rounds of an assay. We call the sequence of fractions a ladder.

\begin{definition}[Rung, ladder, gap]
\label{def:ladder}
A \emph{rung} has a power $p\in[0,1)$. A \emph{ladder} is a finite or
infinite sequence of rungs with powers $(p_i)$. Starting from a gap
$g_0>0$, climbing rung $i$ leaves $g_i=(1-p_i)g_{i-1}$, so
\begin{equation}
g_n=g_0\prod_{i=1}^{n}(1-p_i),\qquad
P_n:=1-\frac{g_n}{g_0}=1-\prod_{i=1}^{n}(1-p_i).
\label{eq:ladder}
\end{equation}
$P_n$ is the \emph{composite power}. A rung with $p\le0$ is refused: it
closes nothing.
\end{definition}

\begin{theorem}[A ladder is a counted sum in logarithmic coordinates]
\label{thm:ladder}
Let $L_n=-\ln(g_n/g_0)=\sum_{i=1}^n-\ln(1-p_i)$. Then:
\begin{enumerate}[label=(\alph*),nosep]
\item each term $-\ln(1-p_i)$ is positive when $p_i>0$, and $L_n$ is
      non-decreasing;
\item no finite ladder closes the gap: $g_n>0$ for every $n$;
\item if $p_i\ge p_{\min}>0$ for all $i$, then $L_n\ge n\ln\frac{1}{1-p_{\min}}$
      and $g_n\to0$;
\item if $\sum_i p_i<\infty$, then $g_n$ converges to a positive limit.
\end{enumerate}
\end{theorem}

\begin{proof}
(a) $0<p_i<1$ gives $0<1-p_i<1$. (b) is a finite product of positive
numbers. (c) applies \cref{prop:floored} to the floored sequence
$-\ln(1-p_i)\ge-\ln(1-p_{\min})$. (d) is the standard criterion for
convergence of infinite products: $\prod(1-p_i)$ with $0\le p_i<1$ converges
to a non-zero limit if and only if $\sum p_i<\infty$~\citep{knopp1990}.
\end{proof}

\Cref{thm:ladder} is \cref{thm:selection} in another guise. A floor on rung
power makes the logarithmic gap a floored sum, which diverges along every
infinite selection of rungs; without a floor, a ladder can climb forever and
stall at a residual. For $p_i=\tfrac12 i^{-2}$ the residual is
$\prod_i(1-\tfrac{1}{2i^2})=\sin(\pi/\sqrt2)/(\pi/\sqrt2)\approx0.3582$; the
numerical product over $500$ rungs gives $0.3585$ (\cref{fig:ladders}B).

\begin{proposition}[A composite does not identify its ladder]
\label{prop:screen}
Ladders with different power sequences can have equal composite powers.
\end{proposition}

\begin{proof}
$(\tfrac12,\tfrac12)$ and $(\tfrac34)$ both give $P=\tfrac34$.
\end{proof}

An observer who records only the end-to-end result of a stepwise process
therefore cannot recover its steps; one who records the gap after each rung
can. In the terms of \cref{sec:model}, the second observer commits once per
rung, the first once per ladder.

\begin{figure*}[t]
\centering
\includegraphics[width=\textwidth]{figures/f4_ladders.pdf}
\caption{\textbf{Ladders.}
(\textbf{A})~Remaining gap after $n$ rungs for a constant ladder, a random
ladder with powers in $[0.05,0.5]$, a harmonic ladder $p_i=0.5/i$ (no floor,
$\sum p_i=\infty$), and a summable ladder $p_i=0.5/i^2$ (no floor,
$\sum p_i<\infty$).
(\textbf{B})~The same ladders as logarithmic gaps $L_n$: floored ladders grow
linearly, the harmonic ladder logarithmically, and the summable ladder stalls
at a residual gap of $0.359$.
(\textbf{C})~Two ladders with equal composite power and different
trajectories. In exact rational arithmetic, none of $200$ random ladders of
$200$ rungs with $p_i\in\{0.01,\dots,0.98\}$ closed the gap.}
\label{fig:ladders}
\end{figure*}

%==============================================================================
\section{Instruments as committers}
\label{sec:devices}
%==============================================================================

Each quantity derived above is read by a familiar kind of instrument, and in
each case the reading is itself a commitment. We list the correspondences to
make clear what an experimental test of each statement would involve; none
of the correspondences is new.

\begin{table*}[t]
\centering
\small
\begin{tabular}{@{}p{3.2cm}p{4.0cm}p{6.4cm}@{}}
\toprule
Quantity & Instrument & What a reading commits\\
\midrule
count $M$ & counter, scaler & one increment per detected event\\
recorded time $t$ & clock & one increment per reference cycle (\cref{sec:time})\\
register entropy & memory, recorder & one symbol per write; resetting a cell
costs at least $\kB T\ln2$ (\cref{sec:entropy})\\
separation cost $\sep$ & separation instruments (chromatography,
mass analysis, electrophoresis) & each resolved fraction is a cut between the
item and its medium; adding an instrument adds contacts
(\cref{thm:monotone})\\
ladder gap $g_n$ & stepwise purification or filtering & the remaining
fraction after each step; the composite alone is a screen
(\cref{prop:screen})\\
\bottomrule
\end{tabular}
\end{table*}

Two consequences are worth stating. First, by \cref{thm:monotone}, adding an
instrument to a measurement setting---adding contacts between the item and
new probes---cannot reduce how well items already in the setting are
separated from their medium; it can only leave the separation unchanged or
raise it. Second, by \cref{cor:no-zeno-records}, a finite instrument resolves
a continuous quantity only to a finite depth; the dichotomy of \cref{sec:zeno}
is one example, and the $\log_2(D/\varepsilon)$ commitments it requires are
the bits of a binary readout of position.

%==============================================================================
\section{Validation}
\label{sec:validation}
%==============================================================================

\Cref{tab:validation} summarises the computational checks. They test the
propositions as stated; they cannot test \cref{post:counted}, which is an
assumption about physical records. All code, the random seeds and the
recorded outputs are provided with the manuscript, and each figure is
regenerated by one command.

\begin{table*}[t]
\centering
\small
\caption{Computational checks. ``Prediction'' is what the cited result
requires; ``observed'' is what the computation returned.}
\label{tab:validation}
\begin{tabular}{@{}p{3.4cm}p{4.6cm}p{5.6cm}@{}}
\toprule
Check & Prediction & Observed\\
\midrule
Counted partial sums ($10^3\times10^4$) & $S_n\ge n$ at every $n$
(\cref{lem:counted}) & $0$ violations; all sequences strictly increasing\\
Running minimum of $U_i/i$ & tends to $0$ (\cref{prop:positivity}) & median
$1.4\times10^{-8}$ after $10^4$ terms\\
Dichotomy & $k(\varepsilon)=\lceil\log_2 D/\varepsilon\rceil$ finite
(\cref{prop:two-series}) & $10,20,\dots,50$ commitments for
$\varepsilon/D=10^{-3},\dots,10^{-15}$\\
$\sep(v)$ under growth & never decreases (\cref{thm:monotone}) & $0$
decreases in $2\times6{,}400$ steps\\
Smallest weight under growth & $1$ (counted); no floor (real)
(\cref{rem:why-counted}) & $1$; median $1.4\times10^{-3}$\\
Kac ring reversal & configuration returns, count does not
(\cref{thm:loschmidt}) & returned at step $300$; register $300$\\
Cat map periods & finite (\cref{prop:recurrence}); $\le3N$~\citep{dyson1992} &
all finite; $\max P/N=3.0$\\
Exact ladders & gap never $0$ (\cref{thm:ladder}b) & $0$ of $200$ closed\\
Summable ladder residual & $\sin(\pi/\sqrt2)/(\pi/\sqrt2)=0.3582$ &
$0.3585$ at $500$ rungs\\
\bottomrule
\end{tabular}
\end{table*}

%==============================================================================
\section{Limitations and scope}
\label{sec:limits}
%==============================================================================

\begin{enumerate}[label=(\arabic*)]
\item \textbf{One physical assumption.} \Cref{post:counted} is assumed, not
derived. Its physical plausibility rests on the thermodynamics of
information~\citep{landauer1961,bennett1982,berut2012}. If records could be
made at arbitrarily small cost, \cref{thm:selection} shows that the
divergence results would fail along suitable selections.
\item \textbf{Records, not dynamics.} The results concern commitment counts,
register entropy and separation costs. They place no constraint on which
dynamical laws hold; in particular they are consistent with time-reversal
invariant dynamics, as \cref{sec:bounded} illustrates.
\item \textbf{The model of individuation is a choice.} Separation cost by
minimum cut is one measure among several (vertex connectivity, spectral
measures, conductance). We chose it because it is computable, monotone under
expansion (\cref{thm:monotone}), and has the max-flow interpretation; other
measures need not share these properties.
\item \textbf{Counted weights are idealisations.} Many physical contact
strengths are real-valued. \Cref{prop:floored} shows that a uniform positive
lower bound suffices; we have not shown that physical contacts have one,
beyond the bound on records in \cref{post:counted}.
\item \textbf{Finite budgets.} \Cref{prop:finite} assumes a finite budget for
distinctions in a part. For a physical system this is a further assumption
about its energy and extent.
\item \textbf{The classical problems.} Our treatment of Zeno and Loschmidt is
a classification of premises by the kind of sum they concern. It is
compatible with, and does not replace, the analytic resolution of the former
and the statistical account of the latter.
\end{enumerate}

%==============================================================================
\section{Conclusion}
\label{sec:conclusion}
%==============================================================================

The content of ``a sum of infinitely many positive integers grows without
bound'' lies in integrality, which supplies a uniform floor, and
\cref{thm:selection} states exactly what the floor guarantees: divergence
along every infinite selection of terms. Modelling records as counted
commitments and keeping them distinct from occupation, we obtained a direction
on records, a recorded time and a register entropy that are three readings of
one count, and an individuation measure on contact graphs that is at least
one and does not decrease as contacts are added. The two classical problems
we examined separate along the same line: Zeno's dichotomy concerns a
convergent measured series and a divergent counted series at once, and
Loschmidt's objection concerns configurations, which recur, rather than
counts, which do not. Ladders of successive refinement inherit the same
structure in logarithmic coordinates.

A companion analysis applies this accounting to a specific molecular system;
the present article is intended to stand on its own and makes no use of
that application.

\paragraph{Data and code availability.}
All computations are in \texttt{validation/} alongside this manuscript
(\texttt{f1\_partial\_sums.py} to \texttt{f4\_ladders.py}); each writes its
numerical record to \texttt{validation/results/} and its figure to
\texttt{figures/}. The computations use NumPy~\citep{harris2020},
SciPy~\citep{virtanen2020}, NetworkX~\citep{hagberg2008} and
Matplotlib~\citep{hunter2007}.

%==============================================================================
\appendix
\section{Protocols of the computations}
\label{app:protocols}
%==============================================================================

\paragraph{F1 (partial sums).} $10^3$ sequences of $10^4$ counted terms
drawn uniformly from $\{1,\dots,5\}$ (seed 20261011). The check counts indices
with $S_n<n$. Real sequences: $c_i=2^{-i}$, $c_i=1/i$ and $c_i=U_i/i$ with
$U_i$ uniform on $(0,1)$; running minima are reported as median and 95\%
interval. The dichotomy uses $D=1$ and $200$ values of $\varepsilon$ from
$10^{-1}$ to $10^{-15}$.

\paragraph{F2 (separation cost).} Each graph starts from four vertices
joined to the medium and to each other with probability $0.6$; at each of
$160$ steps a new vertex is joined to the medium and to three distinct
existing vertices chosen uniformly, and with probability $0.15$ also to the
fixed vertex $v=0$. Counted weights are uniform on $\{1,\dots,5\}$; real
weights are exponential with mean $1.5$. $\sep(v)$ is computed as a maximum
flow (preflow--push) after every step. For the final graphs of ten sequences
per weighting, $\sep$ is computed for every vertex and compared with the
weight of its singleton cut to detect identity regions larger than one
vertex.

\paragraph{F3 (recurrence).} Kac ring with $N=400$ sites, $41$ marked edges
chosen uniformly (seed 3), initial configuration all white. Coarse-grained
entropy $-p\ln p-(1-p)\ln(1-p)$ of the black fraction $p$. Reversal applies
the exact inverse map. The cat map period is the order of
$\left(\begin{smallmatrix}1&1\\1&2\end{smallmatrix}\right)$ modulo $N$, which
is the common period of all orbits.

\paragraph{F4 (ladders).} Four ladders of $500$ rungs as described in
\cref{fig:ladders}. Exact checks use \texttt{fractions.Fraction} with $200$
ladders of $200$ rungs, powers $j/100$, $j$ uniform on $\{1,\dots,98\}$.

\sloppy
\bibliographystyle{unsrtnat}
\bibliography{references}

\end{document}
